\chapter{Population and the Demographic Transition}\label{ch:b3:demographic-transition}

For most of history, more output meant more people, not richer ones. Then, within a
few generations in one country after another, income per person began to rise
steadily, death rates fell, and families became smaller. This chapter sets out the
model that fits the first period, tests it on the data, shows where it stops
working, and asks why fertility fell --- the question on which the evidence is most
divided. It ends with the arithmetic of ageing that the fall in fertility brings.

\begin{openingchart}{GDP per person, United Kingdom series of the Maddison Project,
1252--1850 (2011 dollars).}{Maddison Project Database 2023}
\begin{tikzpicture}
\begin{axis}[omtime, width=0.72\linewidth, height=0.42\linewidth, xmin=1250,
  xmax=1850, ymin=800, ymax=4600, ylabel={Dollars},
  xtick={1300,1400,1500,1600,1700,1800},
  yticklabel style={/pgf/number format/1000 sep={\,}}]
  \addplot[omDef, thick] table[x=year, y=gdppc] {figdata/bachelor-3/33-demographic-transition/out/england.csv};
\end{axis}
\end{tikzpicture}
\end{openingchart}

% figdata: england_stats
The annual series swings by 10 to 15~\% from one year to the next, so the comparisons
below use averages over the ten years around each date. Income per person averaged
1{,}179 dollars around 1300 and 1{,}611 around 1650: 0.09~\% a year over three and a half
centuries. It jumped from 1{,}229 in 1348 to 1{,}519 in 1351, the years of the Black Death, and stayed high for a
century. It then rose by half between 1650 and 1700 (0.85~\% a year), and from the middle of the eighteenth century it grew
in every half-century at more than 0.35~\% a year, reaching 4{,}332 in 1850.

\begin{recall}
Life expectancy, births per woman and ageing are the grades 6--9 volume's; the
Solow model is the Year 1 volume's; overlapping generations, pay-as-you-go pensions
and the old-age dependency ratio are the Year 2 volume's.
\end{recall}

\begin{notation}[Symbols of this chapter]\label{not:b3:demographic-transition:symbols}
$L$ is population; $y$ income per person; $\bar y$ subsistence income; $A$ technology; $n$ the
number of children; $q$ their quality (schooling); $c$ parents' consumption.
\end{notation}

\section{The Malthusian model}\label{sec:b3:demographic-transition:malthus}

\begin{definition}[Malthusian model; Malthusian equilibrium]\label{def:b3:demographic-transition:malthus}
\index{Malthusian model}\index{Malthusian equilibrium}
The \emph{Malthusian model} combines diminishing returns to labour on a fixed area
of land, $y = A L^{-(1 - \alpha)}$ with $0 < \alpha < 1$, and population growth that rises with income,
$\dot L/L = \phi(\ln y - \ln \bar y)$ with $\phi > 0$. Its \emph{Malthusian equilibrium} is the steady state at
which income equals the subsistence level $\bar y$, where population neither grows nor shrinks.
\end{definition}

\begin{proposition}[Malthusian dynamics]\label{prop:b3:demographic-transition:malthus}
With constant $A$, the population converges to $L^* = (A/\bar y)^{1/(1 - \alpha)}$ and income to $\bar y$ from any starting
point. A permanent rise in $A$ raises income at first and population in the long
run, leaving long-run income at $\bar y$.
\end{proposition}

\begin{proof}
$\ln y = \ln A - (1 - \alpha)\ln L$, so $d\ln L/dt = \phi(\ln A - \ln \bar y - (1 - \alpha)\ln L)$, a linear equation in $\ln L$ with coefficient $-\phi(1 - \alpha) < 0$. Its
unique rest point is $\ln L^* = (\ln A - \ln \bar y)/(1 - \alpha)$, and $\ln L - \ln L^*$ decays at rate $\phi(1 - \alpha)$: the steady state is
stable, and there $y = \bar y$. A rise in $A$ raises $y$ at the given $L$; $L$ then grows
until $y$ is back at $\bar y$, with $L^*$ larger by the factor $(A'/A)^{1/(1 - \alpha)}$.
\end{proof}

If technology grows slowly at rate $g$, population grows at $g/(1 - \alpha)$ in the steady state and
income settles just above subsistence, at $\bar y\, e^{g/(\phi(1 - \alpha))}$.

\begin{model}{The Malthusian model}
\assumes Diminishing returns to labour on fixed land; population growth that rises
with income per person; slow technical change.
\predicts Income per person without trend; population tracking technology and land;
higher wages after a fall in population.
\tested Across countries, from 1 to 1500 CE, better technology and more productive land raised
population density but had no significant effect on the standard of living. % ledger: b3dt:ag
In England, builders' real wages were without trend before 1800. % ledger: b3dt:clark
After the Black Death income per person jumped and stayed high for a century (opening
chart).
\breaks Output-based estimates show modest positive growth of income per person in
Britain well before 1800, against the stagnation of real wages; the opening chart's
rise of half in 1650--1700 is the largest such stretch. % ledger: b3dt:bckov
After about 1800 income and population rise together for decades, and then fertility
falls as income rises --- the opposite of the model's second assumption.
\end{model}

The break can be described as three regimes: a Malthusian one, in which population
growth absorbs every gain; a post-Malthusian one, in which population growth absorbs
only part of output growth; and modern growth, after a \omterm{def:b3:demographic-transition:transition}{demographic transition}
reverses the link between income and population growth. % ledger: b3dt:gw

\section{The demographic transition}\label{sec:b3:demographic-transition:transition}

\begin{definition}[Demographic transition]\label{def:b3:demographic-transition:transition}
\index{demographic transition}
The \emph{demographic transition} is the passage from high birth and death rates to
low ones: death rates fall first, population grows quickly for some decades, and
birth rates follow.
\end{definition}

\begin{definition}[Total fertility rate; net reproduction rate]\label{def:b3:demographic-transition:tfr}
\index{total fertility rate}\index{net reproduction rate}
The \emph{total fertility rate} of a year is the sum over ages 15 to 49 of that year's age-specific
birth rates: the number of children a woman would have if she lived through her
childbearing years at those rates. The \emph{net reproduction rate} counts only
daughters and allows for mortality: the number of daughters per woman who survive to
the mothers' ages. A net reproduction rate of 1, a total fertility rate near 2.1 where mortality is low,
replaces the population in the long run.
\end{definition}

\begin{measure}[The total fertility rate]
The rate is a synthetic measure: no woman lives through the rates of a single year.
When women have children later, the rate of each year falls while the number each
woman has over her life may not change, so the period rate understates completed
fertility during a shift to later births. Before 1950 the rates of most countries are
estimates reconstructed from censuses and sparse registers, with the years between
them interpolated; a flat line before the first estimate is a fill-in, not a
measurement of stable fertility.
\end{measure}

\begin{omchart}{Total fertility rate, 1800--2024, six countries (Gapminder to 1949; World
Bank from 1960; the two sources are not joined across the 1950s). Each early series
starts at the source's first estimate --- South Korea 1901, India 1880, Nigeria 1925 ---
because the source repeats that estimate back to 1800; between estimates it
interpolates, so the early lines are smoother than the data behind them.}{Gapminder Fast Track,
children\_per\_woman\_total\_fertility; World Bank, World Development Indicators,
SP.DYN.TFRT.IN}
\begin{tikzpicture}
\begin{axis}[omtime, width=0.72\linewidth, height=0.44\linewidth, xmin=1800, xtick distance=50,
  xmax=2024, ymin=0, ymax=7.5, ylabel={Children per woman}, ytick={0,1,3,4,5,6,7},
  legend style={font=\footnotesize, at={(0.5,-0.2)}, anchor=north}, legend columns=3,
  extra y ticks={2.1}, extra y tick labels={2.1}, extra y tick style={grid=major}]
  \addplot[black, thick] table[x=year, y=tfr] {figdata/bachelor-3/33-demographic-transition/out/tfr_GBR_early.csv};
  \addplot[omDef, thick] table[x=year, y=tfr] {figdata/bachelor-3/33-demographic-transition/out/tfr_FRA_early.csv};
  \addplot[omThm, thick] table[x=year, y=tfr] {figdata/bachelor-3/33-demographic-transition/out/tfr_JPN_early.csv};
  \addplot[omThm, dashed, thick] table[x=year, y=tfr] {figdata/bachelor-3/33-demographic-transition/out/tfr_KOR_early.csv};
  \addplot[omDef, dashed, thick] table[x=year, y=tfr] {figdata/bachelor-3/33-demographic-transition/out/tfr_IND_early.csv};
  \addplot[black!50, dotted, very thick] table[x=year, y=tfr] {figdata/bachelor-3/33-demographic-transition/out/tfr_NGA_early.csv};
  \legend{United Kingdom, France, Japan, South Korea, India, Nigeria}
  \addplot[black, thick, forget plot] table[x=year, y=tfr] {figdata/bachelor-3/33-demographic-transition/out/tfr_GBR_late.csv};
  \addplot[omDef, thick, forget plot] table[x=year, y=tfr] {figdata/bachelor-3/33-demographic-transition/out/tfr_FRA_late.csv};
  \addplot[omThm, thick, forget plot] table[x=year, y=tfr] {figdata/bachelor-3/33-demographic-transition/out/tfr_JPN_late.csv};
  \addplot[omThm, dashed, thick, forget plot] table[x=year, y=tfr] {figdata/bachelor-3/33-demographic-transition/out/tfr_KOR_late.csv};
  \addplot[omDef, dashed, thick, forget plot] table[x=year, y=tfr] {figdata/bachelor-3/33-demographic-transition/out/tfr_IND_late.csv};
  \addplot[black!50, dotted, very thick, forget plot] table[x=year, y=tfr] {figdata/bachelor-3/33-demographic-transition/out/tfr_NGA_late.csv};
\end{axis}
\end{tikzpicture}
\end{omchart}
% figdata: tfr_stats

France's fertility fell earliest, from 4.41 in 1800 to 2.80 in 1900; the United Kingdom's fell from
4.97 to 3.53. South Korea's was 6.0 in 1960 and fell below 2.1 in 1983; in 2024 it was 0.75. India's crossed
2.1 only in 2020; Nigeria's was 4.4 in 2024.

\begin{omchart}{Total fertility rate against GDP per person (PPP, constant 2021 dollars,
logarithmic scale), 2023, 197 countries.}{World Bank, World Development Indicators,
SP.DYN.TFRT.IN and NY.GDP.PCAP.PP.KD}
\begin{tikzpicture}
\begin{axis}[omscatter, width=0.6\linewidth, height=0.4\linewidth,
  xlabel={GDP per person, dollars}, ylabel={Children per woman},
  xtick={3,4,5}, xticklabels={1{,}000, 10{,}000, 100{,}000}, ymin=0, ymax=7]
  \addplot table[x=lx, y=tfr] {figdata/bachelor-3/33-demographic-transition/out/fert_income.csv};
\end{axis}
\end{tikzpicture}
\end{omchart}
% figdata: fert_income_stats

Across countries in 2023, a tenfold difference in income goes with 2.0 fewer children per woman.
The median rate is 4.0 in the poorest fifth and 1.5 in the richest; 54~\% of countries are below 2.1.

\section{Fewer children, more schooling}\label{sec:b3:demographic-transition:quantity-quality}

\begin{definition}[Quantity--quality trade-off]\label{def:b3:demographic-transition:quantity-quality}
\index{quantity--quality trade-off}
The \emph{quantity--quality trade-off} is the link, through the cost of raising
children, between the number of children parents have and what they invest in each:
because the cost of quality is incurred for every child, more children make each
unit of quality dearer, and more quality makes each child dearer.
\end{definition}

\begin{proposition}[Becker--Lewis, Cobb--Douglas case]\label{prop:b3:demographic-transition:becker-lewis}
Parents maximise $\ln c + a\ln n + b\ln q$, with $a > b > 0$, subject to $c + n(\tau + pq) = I$, where $\tau$ is the fixed cost of a
child and $p$ the price of a unit of quality for one child. Then $n = I(a - b)/((1 + a)\tau)$ and $q = b\tau/((a - b)p)$. A rise
in $b$, the weight parents put on quality --- for example because the return to
schooling rises --- lowers $n$ and raises $q$.
\end{proposition}

\begin{proof}
Let $x = n(\tau + pq)$ be spending on children. Given $x$, maximising $a\ln n + b\ln q$ subject to $n(\tau + pq) = x$ gives the conditions
$a/n = \lambda(\tau + pq)$ and $b/q = \lambda np$; dividing, $a pq = b(\tau + pq)$, so $pq = b\tau/(a - b)$, independent of $x$. Then
$a\ln n + b\ln q = a\ln x + \text{constant}$, and maximising $\ln(I - x) + a\ln x$ gives $x = aI/(1 + a)$. So $n = x/(\tau + pq) = x(a - b)/(a\tau) = I(a - b)/((1 + a)\tau)$. Both expressions
are decreasing ($n$) and increasing ($q$) in $b$.
\end{proof}

In this case a rise in income raises the number of children: the model explains the
fall of fertility through the value of schooling, not through income itself. The
question is whether the trade-off exists in the data, and it is hard to test: in
the raw data larger families have less-schooled children, but parents who choose
more children differ in other ways.

\begin{record}[Twins and family size]
Twin births add a child that parents did not plan. In the whole population of
Norway, the negative correlation between family size and education disappears once
birth order is controlled for or twins are used as an instrument; birth order
itself has a large negative effect. % ledger: b3dt:bds
In Israel, with twins and preferences for mixed sexes as instruments, there is no
evidence of a trade-off in schooling, fertility or earnings. % ledger: b3dt:als
\emph{Weakness:} twins differ from other children --- lower birth weight, close
spacing --- which can confound the estimates.
\end{record}

\begin{record}[China's birth policy]
Using twins by birth order and netting out the birth-weight deficit of twins, an
extra child at parity one or two lowers the schooling progress, expected college
enrolment, grades and health of all the children in the family; but the policy's
contribution to China's \omterm{def:b3:human-capital:human-capital}{human capital} was modest. % ledger: b3dt:rz
Using the relaxations of the policy in rural areas, an additional child raised the
first child's school enrolment by about 16 percentage points, more when the children were of the same sex. % ledger: b3dt:qian
The two studies use different sources of variation and find opposite signs; with
the Norwegian and Israeli results, the evidence for a large trade-off at the margin
of family size is weak.
\end{record}

\section{Ageing and the dependency ratio}\label{sec:b3:demographic-transition:ageing}

\begin{definition}[Demographic dividend]\label{def:b3:demographic-transition:dividend}
\index{demographic dividend}
The \emph{demographic dividend} is the rise in income per person that comes from a
rise in the share of the population of working age, as fertility falls and the
large cohorts born before the fall enter work.
\end{definition}

Income per person is output per worker times workers per person: if the working-age
share rises from 50 to 60~\%, income per person rises by 20~\% at unchanged output per worker. The dividend is
paid once: two generations later the same cohorts retire, and the old-age dependency
ratio of the Year 2 volume rises.

\begin{omchart}{Old-age dependency ratio (people aged 65 and over per 100 aged 15--64),
1960--2025, six economies.}{World Bank, World Development Indicators, SP.POP.DPND.OL}
\begin{tikzpicture}
\begin{axis}[omtime, width=0.68\linewidth, height=0.42\linewidth, xmin=1960,
  xmax=2025, ymin=0, ymax=55, ylabel={Per 100 of working age},
  legend pos=north west, legend style={font=\footnotesize}, legend columns=2]
  \addplot[omThm, thick] table[x=year, y=dep] {figdata/bachelor-3/33-demographic-transition/out/dep_JPN.csv};
  \addplot[omDef, thick] table[x=year, y=dep] {figdata/bachelor-3/33-demographic-transition/out/dep_ITA.csv};
  \addplot[omDef, dashed, thick] table[x=year, y=dep] {figdata/bachelor-3/33-demographic-transition/out/dep_FRA.csv};
  \addplot[black, thick] table[x=year, y=dep] {figdata/bachelor-3/33-demographic-transition/out/dep_USA.csv};
  \addplot[omThm, dashed, thick] table[x=year, y=dep] {figdata/bachelor-3/33-demographic-transition/out/dep_CHN.csv};
  \addplot[black!50, dotted, very thick] table[x=year, y=dep] {figdata/bachelor-3/33-demographic-transition/out/dep_NGA.csv};
  \legend{Japan, Italy, France, United States, China, Nigeria}
\end{axis}
\end{tikzpicture}
\end{omchart}
% figdata: dependency_stats

Japan had 9 people of 65 and over per 100 of working age in 1960 and 51 in 2025; China went from 7 to 21; Nigeria,
where fertility is still high, went from 6.0 to 5.4. Under pay-as-you-go pensions, the Year 2
volume's arithmetic applies: the contribution rate equals the replacement rate
times the ratio of pensioners to contributors, so a ratio that rises from 30 to 45 per 100 raises the
rate by half at an unchanged replacement rate.

\section{Exercises}\label{sec:b3:demographic-transition:exercises}

\begin{exercise}[$\star$]\label{exo:b3:demographic-transition:1}
From the opening chart's ten-year averages (1{,}179 dollars around 1300, 1{,}611 around
1650), find the growth of income per person from 1300 to 1650 and say what happens after.
Why would the single years 1300 and 1650 mislead?
\end{exercise}

\begin{exercise}[$\star$]\label{exo:b3:demographic-transition:2}
Income per person is $y = 1{,}000 - 0.5L$ (grain per person, $L$ in thousands) and population grows at $0.0001(y - 400)$ a year. Find
the \omterm{def:b3:demographic-transition:malthus}{Malthusian equilibrium}.
\end{exercise}

\begin{exercise}[$\star$]\label{exo:b3:demographic-transition:3}
There are 40 pensioners per 100 contributors and the replacement rate is 45~\%. Find the pay-as-you-go
contribution rate.
\end{exercise}

\begin{exercise}[$\star$]\label{exo:b3:demographic-transition:4}
What data from before 1800 would contradict the \omterm{def:b3:demographic-transition:malthus}{Malthusian model}?
\end{exercise}

\begin{exercise}[$\star\star$]\label{exo:b3:demographic-transition:5}
In the model of \cref{prop:b3:demographic-transition:malthus} with $\alpha = 1/2$, technology doubles once. Describe the paths of income and
population, and find the long-run change in population.
\end{exercise}

\begin{exercise}[$\star\star$]\label{exo:b3:demographic-transition:6}
Age-specific birth rates per woman per year are 0.02 (15--19), 0.10 (20--24), 0.12 (25--29), 0.09 (30--34), 0.04 (35--39), 0.01 (40--44)
and 0 (45--49). Find the \omterm{def:b3:demographic-transition:tfr}{total fertility rate}. With 48.8~\% of births girls and 98~\% of them surviving to the
mothers' ages, find the \omterm{def:b3:demographic-transition:tfr}{net reproduction rate}.
\end{exercise}

\begin{exercise}[$\star\star$]\label{exo:b3:demographic-transition:7}
Regress the \omterm{def:b3:demographic-transition:tfr}{total fertility rate} on log income across countries in 1960 and in 2023. Compare
the slopes and the means.
\end{exercise}

\begin{exercise}[$\star\star$]\label{exo:b3:demographic-transition:8}
In \cref{prop:b3:demographic-transition:becker-lewis} take $a = 0.5$, and let $b$ rise from 0.2 to 0.3. By what factors do $n$ and $q$ change?
\end{exercise}

\begin{exercise}[$\star\star$]\label{exo:b3:demographic-transition:9}
What result from twin births would contradict the \omterm{def:b3:demographic-transition:quantity-quality}{quantity--quality trade-off}, and
what would support it?
\end{exercise}

\begin{exercise}[$\star\star\star$]\label{exo:b3:demographic-transition:10}
Prove that the \omterm{def:b3:demographic-transition:malthus}{Malthusian equilibrium} is stable when population growth is any
increasing function $h(y)$ with $h(\bar y) = 0$, and find the speed of convergence near it.
\end{exercise}

\begin{exercise}[$\star\star\star$]\label{exo:b3:demographic-transition:11}
Compute the growth of income per person in each half-century from 1300 to 1850 in the
Maddison series, between ten-year averages around each date, and find the first half-century from which growth stays above 0.3~\% a year.
\end{exercise}

\begin{exercise}[$\star\star\star$]\label{exo:b3:demographic-transition:12}
The working-age share rises from 50 to 60~\% of the population over twenty years while output
per worker grows 2~\% a year. Find the growth of income per person, and the share of it due to
the dividend.
\end{exercise}

\section{Problem: Brenland before the Factories}\label{sec:b3:demographic-transition:problem}

\begin{problem}[{Weekend problem --- Brenland before the factories}]\label{pb:b3:demographic-transition:factories}
An invented agrarian Brenland, 1500--1900. Income per person, in kilograms of grain a
year, is $y = A/\sqrt{L}$, with $L$ in thousands. Population grows at $0.1\ln(y/300)$ a year, at most 1.2~\% a year;
from $y = 600$ the maximum falls in a straight line to 0.4~\% at $y = 1{,}200$. Technology grows 0.1~\% a year until 1780
and 1.5~\% after; $A = 9{,}000$ in 1500, with the population at its steady state. A plague in 1600 kills
40~\% of the people. % figdata: agrarian_stats

\medskip\noindent\textbf{Part I --- the Malthusian economy.}
\begin{enumerate}
\item Find the steady-state income while technology grows at 0.1~\% a year.
\item Why is it above 300?
\item Find the population in 1500.
\item Find the population in 1599 and its growth rate.
\item In the long run, what would a doubling of $A$ do to income and population?
\item What does this economy predict for the real wage over a century?
\end{enumerate}

\medskip\noindent\textbf{Part II --- the plague.}
\begin{enumerate}[resume]
\item Find the population just after the plague.
\item Find income in 1600 and the plague's premium over 1599.
\item Check the premium against $y = A/\sqrt{L}$.
\item At what rate does the population grow just after, and why?
\item In what year is half the log premium gone?
\item Find income and population in 1700. Compare with England after 1348.
\end{enumerate}

\medskip\noindent\textbf{Part III --- the take-off.}
\begin{enumerate}[resume]
\item Find income and population in 1780 and 1850.
\item In what year does income first exceed twice subsistence?
\item Find population growth in 1850 and 1899. What has changed?
\item Find income and population in 1900.
\item With $a = 0.5$, $\tau = 0.04$, $p = 0.01$ and income 1, find $n$ and $q$ in \cref{prop:b3:demographic-transition:becker-lewis} for $b = 0.2$.
\item Find them for $b = 0.35$, as the return to schooling rises.
\item Does a rise in income alone lower the number of children in this model?
\end{enumerate}

\medskip\noindent\textbf{Part IV --- ageing in 2025.} Brenland's accounts give employment in 2025
(\code{load("accounts")}); 1{,}950 thousand people are 65 and over, and pensions replace 40~\% of wages.
\begin{enumerate}[resume]
\item Find the number of people of 65 and over per worker.
\item Find the pay-as-you-go contribution rate.
\item If the over-65s grow to 2{,}400 thousand with employment unchanged, find the new rate.
\item What replacement rate would keep the contribution rate unchanged?
\item Why would a rise in employment of older workers help on both sides of the ratio?
\item State the result: the Malthusian steady-state income, the plague's wage premium,
and Brenland's pay-as-you-go contribution rate in 2025.
\end{enumerate}
\end{problem}
